\documentclass[11pt]{article}
\usepackage{worksheet_style}

%\worksheetfilledtrue
\begin{document}
\worksheetheader{24}{Curl and Divergence}

\begin{background}
Green's Theorem involved the quantity $Q_x - P_y$. Today we give it a name: the \textbf{curl}. We also introduce the \textbf{divergence}, which measures expansion or compression. These two differential operators are the vocabulary of the final theorems of the course (Stokes' and the Divergence Theorem).
\end{background}

%=============================================================
\wsection{Curl: Measuring Rotation}

\begin{defbox}{Curl of a Vector Field}
For $\vect{F} = \vec{P, Q, R}$, the \textbf{curl} is (all the same thing, different notation):
\[\text{curl}\,\vect{F} = \curl\vect{F} = \nabla \times \vect{F} = \wbml{\vec{R_y - Q_z,\;\; P_z - R_x,\;\; Q_x - P_y}}\]

\textbf{How to compute} (component by component):
\begin{itemize}[nosep]
\item $\ihat$: $\;\dfrac{\partial R}{\partial y} - \dfrac{\partial Q}{\partial z}$ \qquad $\jhat$: $\;\dfrac{\partial P}{\partial z} - \dfrac{\partial R}{\partial x}$ \qquad $\khat$: $\;\dfrac{\partial Q}{\partial x} - \dfrac{\partial P}{\partial y}$
\end{itemize}

\textit{Memory trick:} each component uses the two variables it doesn't point along, in cyclic order $(x \to y \to z \to x \to \cdots)$. Always subtract in ``reverse'' order. If you know determinants, this equals $\begin{vsmallmatrix}\ihat & \jhat & \khat \\ \partial_x & \partial_y & \partial_z \\ P & Q & R\end{vsmallmatrix}$, but that is not required.

\medskip
-- \textbf{Interpretation:} measures the \wbi{local rotation (``swirling'') of the field at a point}{6cm}

-- $\curl\vect{F} = \vect{0}$ $\implies$ \textbf{irrotational} (no spinning)

-- Direction of $\curl\vect{F}$ = axis of rotation (right-hand rule), magnitude = rate of rotation
\end{defbox}

%=============================================================
\wsubsection{Visualizing Curl}

\begin{center}
\begin{tikzpicture}[scale=0.7, >=Stealth]
  % Positive curl: F = <-y, x>
  \begin{scope}[shift={(-5.5,0)}]
    \draw[->] (-2.3,0) -- (2.3,0); \draw[->] (0,-2.3) -- (0,2.3);
    \foreach \x/\y in {1/0, -1/0, 0/1, 0/-1, 1/1, -1/1, 1/-1, -1/-1, 1.5/0.5, -1.5/-0.5, 0.5/1.5, -0.5/-1.5} {
      \draw[->, bcblue, thick] (\x,\y) -- +(-\y*0.3, \x*0.3);
    }
    \node[below, align=center] at (0,-2.7) {\small\textbf{Positive curl}\\[-1pt]\small $\vec{-y,x}$, curl $= 2$};
  \end{scope}
  % Zero curl: F = <x, y> (source)
  \begin{scope}[shift={(0,0)}]
    \draw[->] (-2.3,0) -- (2.3,0); \draw[->] (0,-2.3) -- (0,2.3);
    \foreach \x/\y in {1/0, -1/0, 0/1, 0/-1, 1/1, -1/1, 1/-1, -1/-1, 1.5/0.5, -1.5/-0.5, 0.5/1.5, -0.5/-1.5} {
      \draw[->, bcgreen!70!black, thick] (\x,\y) -- +(\x*0.3, \y*0.3);
    }
    \node[below, align=center] at (0,-2.7) {\small\textbf{Zero curl}\\[-1pt]\small $\vec{x,y}$, curl $= 0$};
  \end{scope}
  % Negative curl: F = <y, -x>
  \begin{scope}[shift={(5.5,0)}]
    \draw[->] (-2.3,0) -- (2.3,0); \draw[->] (0,-2.3) -- (0,2.3);
    \foreach \x/\y in {1/0, -1/0, 0/1, 0/-1, 1/1, -1/1, 1/-1, -1/-1, 1.5/0.5, -1.5/-0.5, 0.5/1.5, -0.5/-1.5} {
      \draw[->, bcred, thick] (\x,\y) -- +(\y*0.3, -\x*0.3);
    }
    \node[below, align=center] at (0,-2.7) {\small\textbf{Negative curl}\\[-1pt]\small $\vec{y,-x}$, curl $= -2$};
  \end{scope}
\end{tikzpicture}
\end{center}

%=============================================================
\wsubsection{The Paddlewheel Test}

Think of placing a tiny paddlewheel at a point in the flow. The fluid pushes on each blade. If the flow is faster or differently directed on one side than the other, the imbalance creates a torque and the wheel spins.

\begin{itemize}[nosep]
\item \textbf{Positive curl} ($\curl\vect{F} > 0$ in 2D): paddlewheel spins \textbf{CCW}. Right-hand rule: thumb out of page ($+\khat$).
\item \textbf{Negative curl} ($\curl\vect{F} < 0$ in 2D): paddlewheel spins \textbf{CW}. Right-hand rule: thumb into page ($-\khat$).
\item \textbf{Zero curl}: paddlewheel doesn't spin, even if fluid is moving fast.
\end{itemize}

\textbf{Key insight:} curl measures \textit{local spin}, not global shape. A field can have curved streamlines but zero curl, and \textit{straight} streamlines but nonzero curl!

\medskip
\textbf{Example --- shear flow:} In $\vect{F} = \vec{y, 0}$ every arrow is horizontal, yet the paddlewheel spins. The top blade ($y>0$) is pushed right; the bottom blade ($y<0$) is pushed left. This imbalance $\implies$ \textbf{CW} rotation, and indeed $Q_x - P_y = 0 - 1 = -1 < 0$. Negative curl despite no circular streamlines.

\medskip
In 3D, $\curl\vect{F}$ is a vector: its \textbf{direction} is the axis of fastest rotation (right-hand rule) and its \textbf{magnitude} is the rate of rotation.

\begin{example}[Determining Curl Sign from the Paddlewheel]
For each field, predict the paddlewheel behavior at the origin, then verify by computing the 2D curl.

\begin{enumerate}[nosep]
\item[(a)] $\vect{F} = \vec{0,\; x}$
\item[(b)] $\vect{F} = \vec{-y,\; 0}$
\item[(c)] $\vect{F} = \vec{y,\; x}$
\end{enumerate}

\wbbox{
(a) Right of origin ($x>0$): upward push. Left ($x<0$): downward push. $\implies$ \textbf{CCW}.

$Q_x - P_y = 1 - 0 = 1 > 0$ \;\checkmark\; Positive curl (CCW).

\medskip
(b) Above origin ($y>0$): leftward push. Below ($y<0$): rightward push. $\implies$ \textbf{CCW}.

$Q_x - P_y = 0 - (-1) = 1 > 0$ \;\checkmark\; Positive curl (CCW). Different-looking field, same curl as (a)!

\medskip
(c) From $P = y$: top-right, bottom-left (CW tendency). From $Q = x$: right-up, left-down (CCW tendency). They \textbf{cancel}.

$Q_x - P_y = 1 - 1 = 0$ \;\checkmark\; Zero curl. In fact $\vect{F} = \nabla(xy)$ --- conservative.
}{6cm}
\end{example}

%=============================================================
\wsubsection{2D Reduction}

For $\vect{F} = \vec{P, Q}$ (2D, so $R=0$), only the $\khat$ component survives:
$\curl\vect{F} = (Q_x - P_y)\,\khat$. We call the scalar $Q_x - P_y$ the \textbf{2D curl} --- exactly the integrand in Green's Theorem!

%=============================================================
\begin{example}[3D Curl Computation]
Find $\curl\vect{F}$ for $\vect{F} = \vec{xz,\; xyz,\; -y^2}$.

\wbbox{
$P = xz$, \; $Q = xyz$, \; $R = -y^2$.

$\ihat$: $R_y - Q_z = (-y^2)_y - (xyz)_z = -2y - xy$

$\jhat$: $P_z - R_x = (xz)_z - (-y^2)_x = x - 0 = x$

$\khat$: $Q_x - P_y = (xyz)_x - (xz)_y = yz - 0 = yz$

$\curl\vect{F} = \boxed{\vec{-2y - xy,\;\; x,\;\; yz}}$
}{5cm}
\end{example}

\begin{example}[2D Curl: Rotation vs.\ Source]
Compute the 2D curl of (a) $\vect{F} = \vec{-y, x}$ and (b) $\vect{F} = \vec{x, y}$.

\wbbox{
(a) $Q_x - P_y = (x)_x - (-y)_y = 1 - (-1) = 2$. \quad Positive $\implies$ CCW rotation everywhere.

(b) $Q_x - P_y = (y)_x - (x)_y = 0 - 0 = 0$. \quad Irrotational! Despite pointing outward everywhere, there is no local spinning. (This field is conservative: $\vect{F} = \nabla\!\left(\frac{x^2+y^2}{2}\right)$.)
}{3cm}
\end{example}

%=============================================================
\wsection{Divergence: Measuring Expansion}

\begin{defbox}{Divergence of a Vector Field}
For $\vect{F} = \vec{P, Q, R}$, the \textbf{divergence} is (same thing, different notation):
\[\text{div}\,\vect{F} = \div\vect{F} = \nabla \cdot \vect{F} = \wbml{\frac{\partial P}{\partial x} + \frac{\partial Q}{\partial y} + \frac{\partial R}{\partial z}}\]

Just add up how much each component changes along its own direction.

\medskip
-- \textbf{Interpretation:} measures the net \wbi{``outflow'' or expansion at a point}{6cm}

-- $\div\vect{F} > 0$: \textbf{source} (fluid spreading out) \quad $\div\vect{F} < 0$: \textbf{sink} (fluid compressing) \quad $\div\vect{F} = 0$: \textbf{incompressible}
\end{defbox}

%=============================================================
\wsubsection{Visualizing Divergence}

\begin{center}
\begin{tikzpicture}[scale=0.7, >=Stealth]
  % Source: F = <x, y>
  \begin{scope}[shift={(-5.5,0)}]
    \draw[->] (-2.3,0) -- (2.3,0); \draw[->] (0,-2.3) -- (0,2.3);
    \foreach \x/\y in {1/0, -1/0, 0/1, 0/-1, 1/1, -1/1, 1/-1, -1/-1, 1.5/0.5, -1.5/-0.5, 0.5/1.5, -0.5/-1.5} {
      \draw[->, bcblue, thick] (\x,\y) -- +(\x*0.3, \y*0.3);
    }
    \node[below, align=center] at (0,-2.7) {\small\textbf{Source} ($\div > 0$)\\[-1pt]\small $\vec{x,y}$, div $= 2$};
  \end{scope}
  % Sink: F = <-x, -y>
  \begin{scope}[shift={(0,0)}]
    \draw[->] (-2.3,0) -- (2.3,0); \draw[->] (0,-2.3) -- (0,2.3);
    \foreach \x/\y in {1/0, -1/0, 0/1, 0/-1, 1/1, -1/1, 1/-1, -1/-1, 1.5/0.5, -1.5/-0.5, 0.5/1.5, -0.5/-1.5} {
      \draw[->, bcred, thick] (\x,\y) -- +(-\x*0.3, -\y*0.3);
    }
    \node[below, align=center] at (0,-2.7) {\small\textbf{Sink} ($\div < 0$)\\[-1pt]\small $\vec{-x,-y}$, div $= -2$};
  \end{scope}
  % Incompressible: F = <-y, x>
  \begin{scope}[shift={(5.5,0)}]
    \draw[->] (-2.3,0) -- (2.3,0); \draw[->] (0,-2.3) -- (0,2.3);
    \foreach \x/\y in {1/0, -1/0, 0/1, 0/-1, 1/1, -1/1, 1/-1, -1/-1, 1.5/0.5, -1.5/-0.5, 0.5/1.5, -0.5/-1.5} {
      \draw[->, bcgreen!70!black, thick] (\x,\y) -- +(-\y*0.3, \x*0.3);
    }
    \node[below, align=center] at (0,-2.7) {\small\textbf{Incompressible} ($\div = 0$)\\[-1pt]\small $\vec{-y,x}$, div $= 0$};
  \end{scope}
\end{tikzpicture}
\end{center}

\wsubsection{Both Properties / Neither}

Curl and divergence are \textbf{independent}. A field can have any combination:

\begin{center}
\begin{tikzpicture}[scale=0.75, >=Stealth]
  % Both: F = <x-y, x+y> -> div=2, curl=2
  \begin{scope}[shift={(-4.5,0)}]
    \draw[->] (-2.3,0) -- (2.3,0); \draw[->] (0,-2.3) -- (0,2.3);
    \foreach \x/\y in {1/0, -1/0, 0/1, 0/-1, 1/1, -1/1, 1/-1, -1/-1, 1.5/0.5, -1.5/-0.5, 0.5/1.5, -0.5/-1.5} {
      \pgfmathsetmacro{\dx}{(\x-\y)*0.22}
      \pgfmathsetmacro{\dy}{(\x+\y)*0.22}
      \draw[->, violet, thick] (\x,\y) -- +(\dx, \dy);
    }
    \node[below, align=center] at (0,-2.7) {\small\textbf{Both} (curl $\neq 0$, div $\neq 0$)\\[-1pt]\small $\vec{x-y,\; x+y}$\\[-1pt]\small curl $= 2$, div $= 2$};
  \end{scope}
  % Neither: F = <1, 2> uniform flow -> div=0, curl=0
  \begin{scope}[shift={(4.5,0)}]
    \draw[->] (-2.3,0) -- (2.3,0); \draw[->] (0,-2.3) -- (0,2.3);
    \foreach \x/\y in {-1.5/-1.5, -1.5/0, -1.5/1.5, 0/-1.5, 0/0, 0/1.5, 1.5/-1.5, 1.5/0, 1.5/1.5, -0.7/-0.7, 0.7/0.7, -0.7/0.7, 0.7/-0.7} {
      \draw[->, bcgray, thick] (\x,\y) -- +(0.35, 0.7);
    }
    \node[below, align=center] at (0,-2.7) {\small\textbf{Neither} (curl $= 0$, div $= 0$)\\[-1pt]\small $\vec{1,\; 2}$\\[-1pt]\small Uniform flow};
  \end{scope}
\end{tikzpicture}
\end{center}

%=============================================================
\begin{example}[Incompressible Pipe Flow]
A velocity field in a pipe is $\vect{F} = \vec{2x, -y, -z}$. Is the flow incompressible?

\wbbox{
$\div\vect{F} = (2x)_x + (-y)_y + (-z)_z = 2 + (-1) + (-1) = 0$

Yes --- $\div\vect{F} = 0$, so the flow is \textbf{incompressible}. The fluid speeds up along $x$ (pipe narrows) but compensates by slowing in $y$ and $z$. No fluid is created or destroyed.
}{3cm}
\end{example}

\begin{example}[Classifying a Field]
For $\vect{F} = \vec{x, y, z}$, find $\div\vect{F}$ and $\curl\vect{F}$. Interpret.

\wbbox{
$\div\vect{F} = 1 + 1 + 1 = 3$ \quad (source everywhere --- uniform expansion)

$\curl\vect{F}$: \; $\ihat$: $0-0=0$, \; $\jhat$: $0-0=0$, \; $\khat$: $0-0=0$. \; So $\curl\vect{F} = \vect{0}$ (irrotational).

Like an explosion: everything moves outward, nothing rotates.
}{3cm}
\end{example}

\begin{example}[Another Classification]
For $\vect{F} = \vec{-y, x, 0}$, find $\div\vect{F}$ and $\curl\vect{F}$. Interpret.

\wbbox{
$\div\vect{F} = (-y)_x + (x)_y + (0)_z = 0$ \quad (incompressible)

$\curl\vect{F}$: \; $\ihat$: $0-0=0$, \; $\jhat$: $0-0=0$, \; $\khat$: $1-(-1)=2$. \; So $\curl\vect{F} = 2\khat$ (CCW about $z$).

Like a merry-go-round: everything spins, nothing expands or compresses.
}{3cm}
\end{example}

%=============================================================
\wsection{Key Identities}

\begin{formulabox}{Vector Calculus Identities}
\begin{enumerate}[nosep, leftmargin=*]
  \item $\curl(\nabla f) = \wbm{\vect{0}}$ \quad --- conservative fields are always irrotational (by Clairaut's: $f_{xy} = f_{yx}$)

  \textit{Contrapositive:} if $\curl\vect{F} \neq \vect{0}$, then $\vect{F}$ is \textbf{not} conservative.

  \item $\div(\curl\vect{F}) = \wbm{0}$ \quad --- curl fields are always divergence-free (Clairaut's again)

  \item \textbf{Laplacian:} $\nabla^2 f = \div(\nabla f) = \wbm{f_{xx} + f_{yy} + f_{zz}}$ \quad --- divergence of the gradient

  \item \textbf{Product rules:}

  $\nabla(fg) = f\nabla g + g\nabla f$ \qquad $\div(f\vect{F}) = f\,\div\vect{F} + \vect{F}\cdot\nabla f$ \qquad $\curl(f\vect{F}) = f\,\curl\vect{F} + \nabla f \times \vect{F}$
\end{enumerate}
\end{formulabox}

\wsubsection{Operator Relationships}
\[
\text{Scalar field } f \;\xrightarrow{\;\text{grad}\;}\; \text{Vector field } \vect{F} \;\xrightarrow{\;\text{curl}\;}\; \text{Vector field } \vect{G} \;\xrightarrow{\;\text{div}\;}\; \text{Scalar field } h
\]
The two zero identities say: grad $\to$ curl $= \vect{0}$, and curl $\to$ div $= 0$. In other words, the ``derivative of a derivative'' is zero in this chain.

%=============================================================
\begin{example}[Verifying $\curl(\nabla f) = \vect{0}$]
Let $f(x,y,z) = x^2 yz + \sin(yz)$. Verify $\curl(\nabla f) = \vect{0}$.

\wbbox{
\textbf{Gradient:} $\nabla f = \vec{2xyz,\; x^2 z + z\cos(yz),\; x^2 y + y\cos(yz)}$

\textbf{Curl:}

$\ihat$: $(x^2 y + y\cos(yz))_y - (x^2 z + z\cos(yz))_z = [x^2 + \cos(yz) - yz\sin(yz)] - [x^2 + \cos(yz) - yz\sin(yz)] = 0$ \;\checkmark

$\jhat$: $(2xyz)_z - (x^2 y + y\cos(yz))_x = 2xy - 2xy = 0$ \;\checkmark

$\khat$: $(x^2 z + z\cos(yz))_x - (2xyz)_y = 2xz - 2xz = 0$ \;\checkmark

This always works by Clairaut's theorem ($f_{xy} = f_{yx}$, etc.).
}{6cm}
\end{example}

\begin{example}[Proving $\div(\curl\vect{F}) = 0$]
Show that $\div(\curl\vect{F}) = 0$ for any smooth $\vect{F} = \vec{P, Q, R}$.

\wbbox{
$\curl\vect{F} = \vec{R_y - Q_z,\; P_z - R_x,\; Q_x - P_y}$

$\div(\curl\vect{F}) = (R_y - Q_z)_x + (P_z - R_x)_y + (Q_x - P_y)_z$

$= R_{yx} - Q_{zx} + P_{zy} - R_{xy} + Q_{xz} - P_{yz}$

By Clairaut's: $R_{yx} = R_{xy}$, $Q_{zx} = Q_{xz}$, $P_{zy} = P_{yz}$. Every term cancels. $\square$
}{4cm}
\end{example}

%=============================================================
\wsection{Green's Theorem Reformulated}

\begin{formulabox}{Green's Theorem via Curl}
\[\int_a^b \vect{F}(\vect{r}(t))\cdot\vect{r}\,'(t)\,dt = \iint_D (\curl\vect{F})\cdot\khat\,dA\]
The circulation of $\vect{F}$ around $C$ equals the total curl over $D$. This is the 2D special case of \textbf{Stokes' Theorem} (coming soon!).
\end{formulabox}

%=============================================================
\begin{practice}

\prob
Compute $\div\vect{F}$ and $\curl\vect{F}$ for $\vect{F} = \vec{x^2,\, y^2,\, z^2}$.

\wans{$\div\vect{F} = \wbm{2x+2y+2z}$, \quad $\curl\vect{F} = \wbm{\vect{0}}$}

\vspace{1.5cm}

\prob
Show that $\vect{F} = \vec{yz,\, xz,\, xy}$ is both irrotational and incompressible.

\wans{$\curl\vect{F} = \wbm{\vect{0}}$, \quad $\div\vect{F} = \wbm{0}$}

\vspace{1.5cm}

\prob
Find $\curl\vect{F}$ and $\div\vect{F}$ for $\vect{F} = \vec{y^2 z,\; x^2 z,\; x^2 y}$.

\wans{$\curl\vect{F} = \wbml{\vec{0,\; y^2 - 2xy,\; 2z(x-y)}}$, \quad $\div\vect{F} = \wbm{0}$}

\vspace{2.5cm}

\prob
Compute $\curl\vect{F}$ for $\vect{F} = \vec{e^x\cos y,\; -e^x\sin y,\; 0}$. What does the result tell you?

\wans{$\curl\vect{F} = \wbm{\vect{0}}$ \quad $\implies$ \wbm{conservative, potential $f = e^x\cos y$}}

\vspace{2cm}

\prob
Verify $\div(f\vect{F}) = f\,\div\vect{F} + \vect{F}\cdot\nabla f$ for $f = xy$ and $\vect{F} = \vec{1, 1, 0}$.

\wans{both sides $= \wbm{x + y}$ \;\checkmark}

\vspace{1.5cm}

\prob
Prove: if $\vect{F} = \nabla f$ and $f$ is harmonic ($\nabla^2 f = 0$), then $\vect{F}$ is both irrotational and incompressible.

\wbbox{
Irrotational: $\curl\vect{F} = \curl(\nabla f) = \vect{0}$ (identity 1).

Incompressible: $\div\vect{F} = \div(\nabla f) = \nabla^2 f = 0$ (since $f$ is harmonic). $\square$
}{2cm}

\prob
Without computing, predict whether the 2D curl is positive, negative, or zero at the origin. Then verify.

(a) $\vect{F} = \vec{-2y,\; 3x}$ \qquad (b) $\vect{F} = \vec{y,\; y}$ \qquad (c) $\vect{F} = \vec{x^2,\; -y^2}$

\wans{(a) $Q_x - P_y = \wbm{3-(-2) = 5}$ (positive, CCW). \quad (b) $Q_x - P_y = \wbm{0-1 = -1}$ (negative, CW). \quad (c) $Q_x - P_y = \wbm{0}$ (irrotational).}

\vspace{2cm}

\prob
The velocity field of a whirlpool is modeled by $\vect{F} = \vec{-y, x, 2z}$. Find $\curl\vect{F}$ and $\div\vect{F}$. Does this flow rotate, expand, or both?

\wbbox{
$\curl\vect{F} = \vec{0-0,\; 0-0,\; 1-(-1)} = 2\khat$ \quad (CCW rotation about $z$-axis)

$\div\vect{F} = 0 + 0 + 2 = 2$ \quad (expanding vertically)

Both! Like a tornado: spinning while air is pulled upward.
}{3.5cm}

\end{practice}

\end{document}
